\section{A computation described by product schemas}
\label{sec:computation}

We now use a polynomial-size description of an exponentially large
monotone circuit. Its true gates will specify which scalar images must become
singletons in an equivalent graph. We apply the uniform product-circuit
compilation theorem from the companion paper \cite{VariableWL}, with
stronger address padding. The next section
supplies all additional constraints needed for identification.

Fix a binary language in \(\mathrm{EXPTIME}\).  Choose a fixed deterministic
machine \(M\) deciding it, with a single doubly infinite tape. Its input
occupies positions \(0,\ldots,m-1\), all other cells are blank, and its
head starts at zero and moves by at most one cell per step. Its running time is at most
\[
 T=2^{p(m)},
\]
where \(m\) is the input length and \(p\) is a fixed polynomial with integer
coefficients satisfying \(p(m)\ge m+32\).  These conventions cause no loss of
generality.  A fixed multitape machine can be simulated on one tape using
tracks and head markers: sweeps of the used interval find the scanned symbols
and then update the symbols and head positions.  After \(t\) simulated steps,
the interval has length \(O(m+t+1)\), so the overhead is polynomial and can be
absorbed into \(p\).  We also arrange that a halted machine leaves its tape and
head unchanged, retaining its accepting or rejecting state.

For an input word \(w\) of length \(m\), put
\begin{equation}
 d=p(m)+2,\qquad S=2^d=4T,\qquad r=2d.
 \label{eq:computation-parameters}
\end{equation}
In particular $r\ge68$, so the reconstruction and special-twist results
requiring $r\ge32$ apply. Enlarging a valid exponential-time bound to
meet this padding condition changes neither the fixed machine nor its
accepted language.
An address in \(\{0,1\}^r\) consists of two groups of \(d\) bits, in increasing
place order, representing \((\tau,j)\) with \(0\le \tau,j<S\).  Time \(\tau\)
is not cyclic; the position \(j\) is taken modulo \(S\).

We use two kinds of monotone gates.  A scalar gate takes the OR of its incoming
connections and is also true if designated as a seed.  A two-port AND gate
takes the OR of the connections to each port and then the AND of the two port
values.  An empty OR is false.  Each gate type has one gate at every address.
A connection schema specifies its source type, its destination type and port,
a nonempty product source domain, and coordinate injections carrying the
source address to the destination address.  Thus its address data have exactly
the form of a comparison schema from the preceding sections.  A seed or test
request specifies a scalar type and a product of allowed address coordinates.

\begin{proposition}
\label{prop:computation}
For the fixed machine \(M\) and each input \(w\), with parameters
\eqref{eq:computation-parameters}, one can construct in time polynomial in
\(r+m\) an acyclic monotone circuit description with \(O_M(1)\) gate types and
\(O_M((r+1)^2+m)\) connection schemas, seed requests, and test requests.
Every schema has a nonempty product domain and coordinate injections;
the whole description contains \(O_M(r((r+1)^2+m))\) constant-size
coordinate entries. All tests are scalar terminal gates.
The machine \(M\) accepts \(w\) if and only if at least one designated test
gate is true.
\end{proposition}

\begin{proof}
Apply the uniform product-circuit compilation theorem
\cite[Proposition~3.2]{VariableWL} to \(M,p\). It uses the same
one-tape initialization and frozen-halting convention, the same two
least-significant-bit-first coordinate groups, and the same address
length \(r=2(p(m)+2)\). Its position ring has length \(4T=S\),
and its requirement \(p(m)\ge m+2\) is satisfied by our padding.
The theorem gives the stated counts and polynomial-time generation.

The construction uses a fixed local transition rule on three neighboring
tape cells. Scalar gates record cell symbols, and two-port AND gates
recognize the triples producing each next symbol. Time increment and
position increment or decrement are split into carry and borrow cases.
Each case fixes some bits and flips or preserves each coordinate,
so its domain is a product and its coordinate maps are injections.
Combining time and position cases gives \(O_M(r^2)\) connection schemas.
The exceptional input cells and a product decomposition of the blank
suffix supply the seed requests; final-time accepting symbols supply
the terminal tests. The compiler generates these short tables without
listing full addresses. Applied to the decider \(M\) for the chosen
language, its true-test condition is exactly acceptance of \(w\).
\end{proof}
